\section{Introduction}
A Kalman filter (KF) is a type of algorithm which aims to estimate a hidden state in a linear system, given a series of observations. In particular, it handles the case in which the state variable evolves over time, and does so in a mean-squared optimal manner. This project aims to implement such an algorithm in the case of high dimensional linear regression, via the use of principle components analysis (PCA). The combination of PCA and KF exists in the literature, however, the papers largely deal with a one-time PCA fits, rather than updating representations over time. To extend this, we experiment with a new type of rotating covarinace structure, and test a way of keeping the filter in the correct space via a method we term \textit{projection chaining}. We compare projection chaining to the usual static approach and discuss some properties and limitations of the method.